% 07-authorization.tex: проверка не есть полномочия.
\section{Проверка не есть полномочия: состояние, отзыв и размен автономности}
\label{sec:authorization}

\motif{Подлинность удостоверения не есть его нынешние полномочия.}

\analogy{Печать может быть подлинной, а путник всё же нежеланным. Не печать, а доска у ворот говорит, чьи бумаги сегодня в силе, а копии доски, несомой по дороге, верят лишь до тех пор, пока она не устареет.}

\subsection{Два результата, никогда не один}

Весь продукт виден в одном объекте. Доверяющая сторона спрашивает эмитента об удостоверении и
получает для отозванного токена:

\Needspace{9\baselineskip}
\begin{lstlisting}
ПОСТ /апи/в1/проверка  ->  {
  "подлинно": истина,                 // подпись МЛ-ДСА подлинная
  "полномочно\_сейчас": ложь,  // но этот токен отозван
  "состояние": "ОТОЗВАН",
  "пригодно": ложь,                   // подлинный И полномочный -> ложь
  "решение": "отвергнуть"
}
\end{lstlisting}

Подлинность постоянна и поддаётся кэшированию. Полномочия свежи и отзывны. Доверяющей стороне
нужны оба, и ППИ отказывается схлопывать их в один флаг, потому что подпись остаётся подлинной
навсегда, тогда как полномочия под ней могут быть отозваны в любой миг. \sImpl

Конечная точка для доверяющей стороны опознаёт вызывающего как организацию, а не как человека,
через выдачу учётных данных клиента по ОАут2; она принимает доказательство владения, а не
последовательный идентификатор, так что перебрать токены через неё нельзя; она возвращает
единообразный вердикт при несовпадении, так что она не есть оракул существования; и она не хранит
записи о том, кто кого проверял. \sImpl

\subsection{Свежее состояние, а не устаревшая реплика}

Конечная точка проверки направляется на читающую реплику ради пропускной способности, и история
репозитория хранит ту ошибку, которая расщепила два результата навсегда. Подлинность есть свойство
неизменяемого материала, подписанного значения, байтов подписи и хранимого ключа, так что
отстающая реплика не способна ни подделать её, ни отрицать. Пригоден ли токен \emph{сейчас}, есть
вопрос о текущих полномочиях, а отзыв приземляется на ведущий узел; реплика внутри своего окна
отставания могла бы всё ещё показывать только что отозванный токен действующим.
Состояние, решающее \code{пригодно}, закреплено за ведущим узлом, ответ называет свой
источник и несёт \code{по\_состоянию\_на} и нулевую границу отставания, а вызывающий вправе кэшировать
подлинность и обязан перечитывать полномочия. \sImpl

\subsection{Отзыв}

Отзыв есть переход состояния через одну дозволенную процедуру, которая в той же транзакции пишет и
в \entity{СписокОтзывов}, так что список, обращённый к верификатору, никогда не расходится с
состоянием токена, а триггер отвергает всякий иной путь в \code{ОТОЗВАН}. Список несёт канонический
словарь причин; след аудита несёт процедурную подробность и соподписанта, когда потолок усмотрения
эмитента его потребовал. Отзыв смотрит вперёд: события, записанные, пока удостоверение было
действительным, не обессмысливаются, потому что история не переписывается. \sImpl

\subsection{Полномочия при выключенной сети}

\figfloat{offline-status}{Размен автономности, нарисованный как временная шкала. Эмитент
подписывает утверждение о состоянии с окном; держатель подшивает его к предъявлению; доверяющая
сторона, проверяющая автономно внутри окна, принимает. Если эмитент отзовёт удостоверение внутри
окна, последнее утверждение \code{ДЕЙСТВУЮЩИЙ} останется действительным до истечения, и автономный
верификатор об этом знать не может. Окно ограничено сроком жизни, выданным эмитентом (по умолчанию
час), и величиной \code{макс\_окно}, которую навязывает верификатор; доверяющая сторона с высоким
риском требует секунд либо выходит в сеть. Кто выбирает окно и кто его несёт: эмитент задаёт срок
жизни, а каждая доверяющая сторона свой \code{макс\_окно}; вред долгого окна ложится на того, кто
полагается на удостоверение, отозванное внутри окна, и на держателя, которого отзыв должен был
защитить. Полярис делает окно явным и ограниченным; какая длина приемлема, есть решение управления,
которого он не принимает.}{fig:offline-status}

Автономный верификатор не способен спросить эмитента. Полярис отвечает на это
короткоживущим подписанным утверждением о состоянии: подпись эмитента тем же ключом, которым
подписано удостоверение, над текущим состоянием токена и окном. Держатель забирает его,
предъявляя подлинное удостоверение, без предъявителя и без личных данных, и подшивает к
предъявлению. Верификатор принимает автономно, только если удостоверение подлинно, утверждение
подлинно, связано именно с этим удостоверением, свежо по собственным часам верификатора, не шире
окна, которое верификатор примет, и равно \code{ДЕЙСТВУЮЩИЙ}. Правило, что верификатор никогда не
принимает удостоверения позже чем через одно окно после его отзыва и что окном является то,
которое задал верификатор, а не то, которое вычеканил эмитент, есть теперь формальная модель,
проверяемая при каждой отправке с конфигурацией, на которой она обязана упасть. \sImpl

Цена изложена так же прямо, как и выгода. Поскольку проверка не касается эмитента, эмитент никогда
не узнаёт, что она была: автономно приватнее, чем в сети. И поскольку она не касается эмитента,
последнее утверждение отозванного удостоверения остаётся действительным до истечения.
Ненаблюдаемость эмитентом, автономная доступность и свежесть отзыва не могут быть доведены до
предела одновременно; развёртывание выбирает две, а окно есть число политики, а не закон. Одно
следствие записано в сводной таблице как известное ограничение, а не исправлено: эмитент чеканит
часовое окно, собственная граница по умолчанию у эталонного верификатора не задана, а потолок,
который несёт сам предмет, проверкой утверждения не читается, так что предельный автономный срок
жизни состояния держится отчасти на настройках доверяющей стороны. Доверяющая сторона, которой
нужны секунды, задаёт секунды либо выходит в сеть.

\subsection{Состояние между органами и через сети, которым не доверяют}

Ещё два подписанных объекта переносят состояние между органами без обратного вызова, который
эмитент мог бы записать. Контрольная точка эпохи есть обязательство органа о последней точке его
цепи эпох: номер эпохи, её корень Меркла и предыдущая эпоха, которую она продолжает; две
контрольные точки доказывают монотонность, а два разных корня, подписанные под одним номером
эпохи, суть криптографическое доказательство того, что орган двоесказательствовал о собственной
истории. Лента отзывов есть упорядоченное множество листьев отозванных удостоверений, каждый из
которых есть ША3-256 значения токена, выводимый лишь держателем; более новая лента, выбросившая
опубликованный отзыв, есть пойманный откат, потому что лежащая под ней таблица работает только на
добавление. В масштабе федерации сводчик зеркалит ленты и контрольные точки многих органов,
дословно и по-прежнему под собственной подписью каждого участника, в один короткоживущий набор.
Сводчику не доверяют в части правильности: полностью распоряжаясь набором, пересчитывая
обязательство и переподписывая конверт, он всё равно не способен подделать состояние участника,
потому что не способен подписаться как участник. Набор добавляет доступность, а не доверие.
\sImpl

Подписание предмета о состоянии есть также и то, что позволяет недоверенному кэшу его переносить,
а кэш вносит тот единственный отказ, которого подписание не предотвращает, а именно время.
Директива кэширования никогда не есть константа: \code{макс-возраст} всякого оконного
предмета есть его собственный остаток жизни, вычисленный из подписанного эмитентом
\code{истекает\_в}, так что кэш не способен пережить окно, к которому эмитент себя обязал; всякий
предмет, называющий одно удостоверение, есть \code{не-хранить}, потому что общий кэш, держащий
такой, выдал бы удостоверение одного держателя другому; а истёкшее или неразбираемое окно не
кэшируется вовсе. Раздача в масштабе сети доставки содержимого тем самым описана и отработана в
учениях; ни в одном развёртывании её нет. \sImplMark{} для предметов, учений и правил кэширования;
\sFut{} для развёртывания раздачи.

Темп смены эпох есть то единственное число, которое задаёт три не связанных между собой свойства,
и проектная запись теперь говорит, какие именно: свежесть отзыва, поскольку отозванное
удостоверение выпадает из \emph{следующей} эпохи; размер толпы, поскольку множество анонимности
есть эпоха; и то, как часто обнуляется у доверяющей стороны учёт <<один человек однажды>>.
Закрытие эпохи было когда-то дорогой частью, и дорогой по дурной причине: дерево постоянной
глубины, дополненное до национальной глубины 24, материализовывало шестнадцать миллионов листьев
при любом населении, так что корень над тысячей участников занимал 10,8 секунды и 2,9~ГБ.
Дополнение есть одно повторяющееся значение, хешируемое по разу на уровень, и
корень над миллионом участников занимает 2,1 секунды в десятках мегабайт; связывающим ограничением
на темп есть нижний предел анонимности на контекст, а не вычисления. \sImpl

\begin{layercard}{Карточка слоя: проверка и состояние}
\qa{Что это}{Конечная точка проверки для доверяющей стороны, список отзывов, подписанное утверждение о состоянии, контрольные точки эпох, ленты отзывов, сводный набор состояний и правила кэширования, позволяющие недоверенной сети всё это переносить.}
\qa{Зачем оно существует}{Подлинность не есть полномочия. Верификатору нужно знать не только, что утверждение было сделано, но и что оно по-прежнему в силе, а иногда нужно знать это при выключенной сети.}
\qa{Чему доверяет}{Подписывающему ключу эмитента, ради утверждения; собственным часам, ради свежести; ведущей базе данных, ради сетевого вердикта.}
\qa{Чему отказывает}{Мнению реплики о текущем состоянии; утверждению шире потолка верификатора; слову сводчика о состоянии участника; ленте, не связанной с ключом эмитента; записи кэша, которая пережила бы подписанное эмитентом окно.}
\qa{Что видит}{Состояние одного удостоверения в один момент; множество отозванных листьев.}
\qa{Чего не видит}{Отзыв, случившийся после того, как было подписано имеющееся у него утверждение; действующее население, которого не публикует ни одна лента.}
\qa{Если откажет}{Устаревший вердикт ограничен окном, и граница есть граница верификатора; раздвоенная цепь эпох или откаченная лента суть неотрицаемое свидетельство против органа; лгущий сводчик ловится на подписи участника.}
\qa{Сейчас или потом}{Сейчас, для предметов и правил. Раздача в национальном масштабе потом.}
\qa{Внутри и снаружи}{Внутри: верификаторы. Снаружи: граница приватности, решающая, что верификатору дозволено узнать, пока он проверяет.}
\end{layercard}

\nextlayer{Состояние доказывает нынешние полномочия, но доказательство способно открыть куда
больше, чем верификатору требовалось, и всякая проверка есть факт о чьём-то дне. Следующий слой
ограничивает то, что узнаёт кто бы то ни было, и настоящая версия сообщает, что об этой границе
было измерено, а не что задумывалось.}
